\chapter{Trajetória e compatibilidade dos materiais}\label{ap:historico}
A monografia consolida o trabalho existente, mas não trata suas fases como uma sequência de confirmações independentes. A mudança de universo altera o significado das métricas e impede reunir resultados antigos e recentes em uma única comparação de desempenho.

\begin{table}[htbp]\centering\small
\caption{Papéis das etapas do projeto na versão consolidada.}\label{tab:historico}
\begin{tabularx}{\linewidth}{>{\raggedright\arraybackslash}p{0.24\linewidth}>{\raggedright\arraybackslash}p{0.22\linewidth}Y}\toprule
Material & Universo ou unidade & Uso nesta monografia \\\midrule
Integração inicial Spotify--LFM2b & 111.758 pares aceitos & Contexto de origem, separação de respostas e princípios de matching; sem métricas misturadas à comparação recente.\\
Desenvolvimento e teste v0.5/v0.6 & 7.648 e 2.592 observações & História de exposição do teste e decisões; não confirmação independente da coorte atual.\\
Diagnósticos v0.7--v0.9 & Recortes e ajustes históricos & Registro de exploração posterior ao teste; gráficos antigos não representam erros atuais.\\
Coorte aceita antiga ampliada & 58.633 observações & Pacote parcial preservado, de outro universo e outra combinação de variáveis.\\
Amostra recente e recuperação v2 & 200.000 / 77.535 faixas & Linha principal: descrição, seleção textual e quatro configurações Ridge comuns.\\
Diagnósticos e parecer recentes & Previsões externas existentes & Complementos condicionais, resíduos, clipping e auditoria do prefixo, sem novos treinos.\\
\bottomrule\end{tabularx}
\fonte{Relatórios e resumos históricos do próprio projeto.}
\end{table}

Foram reaproveitados fundamentos de preservação dos originais, controle de identidade, separação de transformações por treino, comparação pareada e transparência da proveniência. Esses princípios permaneceram compatíveis. As definições de população, textos, vetores e blocos antigos não foram transportadas por simples igualdade de ID ou dimensão.

Os dois pacotes anteriores do Overleaf eram trechos parciais de métodos, resultados e discussão. Sua estrutura editorial foi considerada na organização, mas seus universos não foram fundidos. A versão consolidada incorpora os resultados atuais e o parecer posterior à leitura do PDF, com capítulos introdutórios, fundamentação, fontes, seleção, discussão e apêndices acrescentados.

O inventário de reaproveitamento identifica os arquivos acadêmicos existentes e classifica sua compatibilidade. O acervo histórico foi mantido, inclusive materiais que não entram como evidência principal. A preservação torna possível reconstruir a trajetória sem reinterpretar retrospectivamente os experimentos como se tivessem sido planejados todos ao mesmo tempo.
